\chapter{动态规划：状态、转移和重复子问题}

动态规划不是“看到最值就开数组”。它解决的是有重叠子问题、且大问题能由小问题稳定推出的问题。暴力搜索之所以慢，是因为同一个子问题被不同路径反复计算；动态规划的优化，就是给子问题命名、保存结果，并按照依赖顺序把结果推出来。学 DP 最重要的五件事是：状态表示什么，答案在哪里，初始状态是什么，转移从哪里来，遍历顺序为什么满足依赖。

爬楼梯是最小的 DP 模型。题目问到达第 \code{n} 阶有多少种方法，每次可以爬 1 或 2 阶。暴力递归会写成 \code{ways(n) = ways(n - 1) + ways(n - 2)}，但 \code{ways(n - 2)}、\code{ways(n - 3)} 等子问题会被反复计算，形成指数级时间。优化方法是从小到大保存每个阶数的答案。

\begin{lstlisting}[language=C++,caption={爬楼梯：滚动变量保存最近两个状态}]
int climb_stairs(int n)
{
    if (n <= 2) {
        return n;
    }

    int previous_two = 1;
    int previous_one = 2;
    for (int step = 3; step <= n; ++step) {
        const int current = previous_one + previous_two;
        previous_two = previous_one;
        previous_one = current;
    }

    return previous_one;
}
\end{lstlisting}

这里状态可以写成 \code{dp[i]}：到达第 \code{i} 阶的方法数。最后一步要么从 \code{i - 1} 走 1 阶，要么从 \code{i - 2} 走 2 阶，因此转移是 \code{dp[i] = dp[i - 1] + dp[i - 2]}。由于每个状态只依赖前两个状态，可以把数组压成两个变量。时间复杂度 \complexity{O(n)}，空间复杂度 \complexity{O(1)}。这题的价值不在难度，而在于展示 DP 的完整骨架。

打家劫舍把“相邻不能同时选”的约束放进状态。暴力方法是对每间房决定偷或不偷，枚举 \complexity{2^n} 个子集，再过滤相邻冲突。更好的建模是让 \code{dp[i]} 表示考虑前 \code{i} 间房能偷到的最大金额。第 \code{i} 间房只有两种结局：不偷它，答案是 \code{dp[i - 1]}；偷它，就不能偷第 \code{i - 1} 间，答案是 \code{dp[i - 2] + nums[i - 1]}。

\begin{lstlisting}[language=C++,caption={打家劫舍：选与不选的状态转移}]
int rob(const std::vector<int>& nums)
{
    int previous_two = 0;
    int previous_one = 0;

    for (const int money : nums) {
        const int current = std::max(previous_one, previous_two + money);
        previous_two = previous_one;
        previous_one = current;
    }

    return previous_one;
}
\end{lstlisting}

这段代码把 \code{dp[i - 2]} 和 \code{dp[i - 1]} 压成两个变量。遍历到当前房屋时，\code{previous_one} 是不偷当前房屋时可继承的最优值，\code{previous_two + money} 是偷当前房屋时的最优值。正确性来自最优子结构：当前决策之后，剩余可用的历史范围只与前一间或前两间有关，不需要知道更早的具体选择。易错点是下标：如果用数组，建议定义 \code{dp[0] = 0}，\code{dp[1] = nums[0]}，这样第 \code{i} 个状态表示前 \code{i} 间房，避免把房屋下标和状态下标混在一起。

零钱兑换展示“最少数量”型 DP。题目给硬币面额和目标金额，要求凑成目标金额所需的最少硬币数。暴力搜索会尝试所有硬币序列，许多不同顺序会到达同一个剩余金额，例如先选 1 再选 2，与先选 2 再选 1，都可能进入同一个子问题。状态应定义为 \code{dp[amount]}：凑出金额 \code{amount} 的最少硬币数。

\begin{lstlisting}[language=C++,caption={零钱兑换：从可达金额推出新金额}]
int coin_change(const std::vector<int>& coins, int amount)
{
    constexpr int INF = 1'000'000'000;
    std::vector<int> dp(amount + 1, INF);
    dp[0] = 0;

    for (int value = 1; value <= amount; ++value) {
        for (const int coin : coins) {
            if (coin > value || dp[value - coin] == INF) {
                continue;
            }
            dp[value] = std::min(dp[value], dp[value - coin] + 1);
        }
    }

    return dp[amount] == INF ? -1 : dp[amount];
}
\end{lstlisting}

转移的含义是：如果最后一枚硬币是 \code{coin}，那么前面必须先凑出 \code{value - coin}，所以候选答案是 \code{dp[value - coin] + 1}。这里用一个足够大的 \code{INF} 表示不可达，而不是用 \code{-1} 参与最小值计算，因为 \code{-1 + 1} 会污染转移。时间复杂度 \complexity{O(amount \cdot c)}，\code{c} 是硬币种类数。完全背包类题还要区分“组合数”和“排列数”：本题只问最少硬币数，内外循环顺序不会影响最优值；如果问方案数量，循环顺序就会改变计数语义。

最长递增子序列有两种经典解法。第一种 DP 定义 \code{dp[i]} 为以 \code{nums[i]} 结尾的最长递增子序列长度。为了求它，需要看所有 \code{j < i} 且 \code{nums[j] < nums[i]} 的位置，转移为 \code{dp[i] = max(dp[j] + 1)}。这比枚举所有子序列的 \complexity{2^n} 好很多，但仍是 \complexity{O(n^2)}。

\begin{lstlisting}[language=C++,caption={最长递增子序列：二次 DP}]
int length_of_lis_dp(const std::vector<int>& nums)
{
    if (nums.empty()) {
        return 0;
    }

    std::vector<int> dp(nums.size(), 1);
    int answer = 1;

    for (int i = 0; i < static_cast<int>(nums.size()); ++i) {
        for (int j = 0; j < i; ++j) {
            if (nums[j] < nums[i]) {
                dp[i] = std::max(dp[i], dp[j] + 1);
            }
        }
        answer = std::max(answer, dp[i]);
    }

    return answer;
}
\end{lstlisting}

第二种优化来自“只关心长度对应的最小结尾”。维护数组 \code{tails}，其中 \code{tails[len - 1]} 表示长度为 \code{len} 的递增子序列中，最小可能结尾值。结尾越小，后面越容易接上更大的数。遍历新数 \code{x} 时，用二分找到第一个大于等于 \code{x} 的位置并替换；如果找不到，说明 \code{x} 可以扩展出更长序列。

\begin{lstlisting}[language=C++,caption={最长递增子序列：贪心尾值加二分}]
int length_of_lis_binary_search(const std::vector<int>& nums)
{
    std::vector<int> tails;

    for (const int x : nums) {
        auto position = std::lower_bound(tails.begin(), tails.end(), x);
        if (position == tails.end()) {
            tails.push_back(x);
        } else {
            *position = x;
        }
    }

    return static_cast<int>(tails.size());
}
\end{lstlisting}

这个算法经常被误解成 \code{tails} 本身就是最终的子序列。它不一定是。它保存的是每个长度下最有潜力的最小结尾。替换操作不会减少已经找到的长度，只会让同长度的结尾更小，从而给未来元素更多机会。时间复杂度 \complexity{O(n\log n)}，空间复杂度 \complexity{O(n)}。如果面试官追问为什么能替换，要回答：相同长度下，更小结尾不会比更大结尾更差，因为它能接纳的后续元素集合只会更多。

单词拆分把字符串前缀变成状态。题目给字符串和词典，问字符串能否被切成若干词典单词。暴力搜索会尝试每个切分点，遇到相同后缀时重复判断。状态 \code{dp[i]} 表示前 \code{i} 个字符是否可以被合法拆分。若存在 \code{j < i}，使得 \code{dp[j]} 为真且子串 \code{s[j..i)} 在词典中，则 \code{dp[i]} 为真。

\begin{lstlisting}[language=C++,caption={单词拆分：前缀可达状态}]
bool word_break(const std::string& s, const std::vector<std::string>& word_dict)
{
    std::unordered_set<std::string> words(word_dict.begin(), word_dict.end());
    std::vector<char> dp(s.size() + 1, false);
    dp[0] = true;

    for (std::size_t right = 1; right <= s.size(); ++right) {
        for (std::size_t left = 0; left < right; ++left) {
            if (!dp[left]) {
                continue;
            }
            if (words.contains(s.substr(left, right - left))) {
                dp[right] = true;
                break;
            }
        }
    }

    return dp[s.size()];
}
\end{lstlisting}

这里 \code{dp[0] = true} 代表空前缀可拆分，是所有从 0 开始的单词能够成立的基础。代码直接使用 \code{substr}，写法清晰，但会复制子串；如果输入规模更大，可以限制最大单词长度，或者用字符串视图和自定义哈希、字典树来减少复制。复杂度通常写成 \complexity{O(n^3)} 的保守形式，因为有两层切分枚举，子串构造和哈希也与长度有关；若把子串查询视作均摊常数，则是 \complexity{O(n^2)}。面试中最好主动说明这个实现的复制成本。

最长公共子序列是二维 DP 的代表。子序列可以删除字符但不能改变相对顺序。暴力方法会枚举两个字符串的所有子序列再比较，数量指数级。状态 \code{dp[i][j]} 表示 \code{text1} 的前 \code{i} 个字符和 \code{text2} 的前 \code{j} 个字符的最长公共子序列长度。

\begin{lstlisting}[language=C++,caption={最长公共子序列：二维前缀 DP}]
int longest_common_subsequence(const std::string& text1,
                               const std::string& text2)
{
    const int rows = static_cast<int>(text1.size());
    const int cols = static_cast<int>(text2.size());
    std::vector<std::vector<int>> dp(rows + 1, std::vector<int>(cols + 1, 0));

    for (int i = 1; i <= rows; ++i) {
        for (int j = 1; j <= cols; ++j) {
            if (text1[i - 1] == text2[j - 1]) {
                dp[i][j] = dp[i - 1][j - 1] + 1;
            } else {
                dp[i][j] = std::max(dp[i - 1][j], dp[i][j - 1]);
            }
        }
    }

    return dp[rows][cols];
}
\end{lstlisting}

为什么相等时来自左上角？因为当前两个字符可以配成公共子序列的最后一个字符，剩下问题变成两个前缀各去掉这个字符。为什么不等时取上方和左方最大值？因为两个末尾字符不能同时作为同一个结尾，至少要丢掉其中一个。多出来的第 0 行和第 0 列代表空串，避免了边界分支。时间复杂度 \complexity{O(mn)}，空间复杂度 \complexity{O(mn)}；如果只要长度，可以滚动数组压到 \complexity{O(n)}。

编辑距离比 LCS 更综合。题目允许插入、删除、替换，把一个单词变成另一个单词，求最少操作数。状态 \code{dp[i][j]} 表示把 \code{word1} 的前 \code{i} 个字符变成 \code{word2} 的前 \code{j} 个字符的最少操作数。这个定义必须精确，因为三种操作都围绕“前缀末尾”展开。

\begin{lstlisting}[language=C++,caption={编辑距离：插入、删除、替换三种转移}]
int min_distance(const std::string& word1, const std::string& word2)
{
    const int rows = static_cast<int>(word1.size());
    const int cols = static_cast<int>(word2.size());
    std::vector<std::vector<int>> dp(rows + 1, std::vector<int>(cols + 1, 0));

    for (int i = 0; i <= rows; ++i) {
        dp[i][0] = i;
    }
    for (int j = 0; j <= cols; ++j) {
        dp[0][j] = j;
    }

    for (int i = 1; i <= rows; ++i) {
        for (int j = 1; j <= cols; ++j) {
            if (word1[i - 1] == word2[j - 1]) {
                dp[i][j] = dp[i - 1][j - 1];
                continue;
            }

            const int delete_cost = dp[i - 1][j] + 1;
            const int insert_cost = dp[i][j - 1] + 1;
            const int replace_cost = dp[i - 1][j - 1] + 1;
            dp[i][j] = std::min({delete_cost, insert_cost, replace_cost});
        }
    }

    return dp[rows][cols];
}
\end{lstlisting}

初始化含义是：把长度为 \code{i} 的前缀变成空串，需要删除 \code{i} 次；把空串变成长度为 \code{j} 的前缀，需要插入 \code{j} 次。若末尾字符相同，不需要新增操作，直接继承左上角；若不同，删除对应上方，插入对应左方，替换对应左上角。时间和空间都是 \complexity{O(mn)}。易错点是从哪个字符串视角解释插入和删除：只要状态定义固定，\code{dp[i - 1][j] + 1} 就是删掉 \code{word1} 的末尾，\code{dp[i][j - 1] + 1} 就是在 \code{word1} 末尾插入 \code{word2[j - 1]}。

动态规划题经常可以做空间优化，但不要为了炫技先写滚动数组。二维表能展示依赖关系，滚动数组则要求你非常清楚哪些旧值还要用、哪些可以覆盖。面试时，先给出清晰状态和二维实现，再说明可压缩空间，通常比一上来写难读的一维代码更稳。

\begin{principlebox}{动态规划题复盘模板}
每道 DP 题按五步复盘：第一，状态是什么，必须能用自然语言解释；第二，答案在 \code{dp[n]}、\code{dp[m][n]} 还是某个最大值里；第三，空状态和最小状态如何初始化；第四，转移枚举的是最后一步、最后一个选择，还是最后一段区间；第五，遍历顺序是否保证依赖状态已经算完。若你只能背出公式，却说不出公式对应的最后一步选择，这题就还没有真正掌握。
\end{principlebox}
